\section{历史损益、现金与营运资本：先看周期，再看增长}

2023--2025 年的报表说明，雅化的收入、利润和经营现金并不同步。2024 年收入下滑后归母亏损，2025 年收入温和修复、归母回正，但经营现金转为 -5.70 亿元。任何把 2026H1 预告直接外推到多年高增长的估值，都必须先解释这条利润--现金差异。

\begin{exhibitbox}[表 A-2：历史财务摘要]
\begin{tabularx}{\exhibitboxwidth}{L{3.3cm}R{2.1cm}R{2.1cm}R{2.1cm}R{2.1cm}}
\toprule
项目（亿元，除 EPS） & 2023A & 2024A & 2025A & 2026Q1 \\
\midrule
营业收入 & 118.95 & 77.16 & 85.43 & 28.30 \\
归母净利 & 0.40 & -2.57 & 6.32 & 3.39 \\
经营现金流 & 8.31 & 9.44 & -5.70 & -4.31 \\
期末应收账款 & -- & -- & 13.26 & 19.80 \\
期末存货 & -- & -- & 16.96 & 18.73 \\
归母权益 & -- & -- & 106.95 & 110.33 \\
\bottomrule
\end{tabularx}
\end{exhibitbox}

\section{2026Q1 资产负债表：净现金代理不等于自由现金流}

2026Q1 货币资金加交易性金融资产 26.03 亿元；短期借款、一年内到期非流动负债和长期借款合计 11.19 亿元，形成约 14.84 亿元净现金代理。该数字被用于 EV/EBITDA 的交叉检查，但不作为目标价的独立加项：首先它不等于所有资金可自由分配；其次 Q1 的营运资本占用正在侵蚀现金创造；最后资源、项目与扩产对资本的实际需求尚待正式披露。

\begin{exhibitbox}[表 A-3：2026Q1 资本与营运资本观察]
\begin{tabularx}{\exhibitboxwidth}{L{4.2cm}R{2.4cm}X}
\toprule
项目 & 数值 & 研究含义 \\
\midrule
总资产 & 150.75 亿元 & 资产规模不等于可直接估值资产 \\
总负债 & 37.65 亿元 & 需结合期限与营运资本分析 \\
货币资金 + 交易性金融资产 & 26.03 亿元 & EV 交叉检查的现金代理 \\
有息负债代理 & 11.19 亿元 & 短借 + 一年内到期 + 长借 \\
净现金代理 & 14.84 亿元 & 不单独加到股权目标价 \\
应收账款较年初变动 & +6.54 亿元 & H1 必须检查回款与收入匹配 \\
存货较年初变动 & +1.77 亿元 & 需辨识价格、备货与减值风险 \\
\bottomrule
\end{tabularx}
\end{exhibitbox}

\section{历史分部：用已实现利润池约束前瞻假设}

2025 年锂产品收入 48.24 亿元、毛利 5.89 亿元；民爆产品收入 21.41 亿元、毛利 10.11 亿元；爆破服务收入 11.75 亿元、毛利 2.31 亿元。民爆产品和服务合计毛利 12.42 亿元，是模型中非锂利润池的基础，但公司未披露相应的分部归母、现金流和项目回款，故不将其拆成独立 SOTP 资产。运输毛利率 13.29\% 为合并成本派生代理，应仅作方向性使用。

\sourcenote{公司 2023--2025 年年报与 2026Q1 报告；结构化财务提取与来源页码见案例数据包。}
